% Part: computability
% Chapter: computability-theory
% Section: fixed-point-thm

\documentclass[../../../include/open-logic-section]{subfiles}

\begin{document}

\olfileid{cmp}{thy}{fix}
\olsection{స్థిరబిందు సిద్ధాంతం}

ఆగే సమస్యను మళ్లీ పరిశీలిద్దాం. తాత్కాలిక సంకేతంగా,
$\tuple{x,y}$కు బదులు $\gn{\cfind{x}(y)}$ అని రాద్దాం; దీన్ని
$\cfind{x}(y)$ విలువకు ఒక ``పేరు''గా భావించండి. ఈ సంకేతంతో, ఆగే
సమస్య అనిర్ణయనీయమని మనం ఇచ్చిన నిరూపణల్లో ఒకదాన్ని తిరిగి
వ్యక్తపరచవచ్చు.

ప్రశ్న: ప్రతి $x$, $y$లకు కింది ధర్మం గల గణనీయ ప్రమేయం~$h$ ఉందా?
\[
h(\gn{\cfind{x}(y)}) =
\begin{cases}
1 & \text{$\cfind{x}(y) \fdefined$ అయితే} \\
0 & \text{ఇతర సందర్భాల్లో.}
\end{cases}
\]

సమాధానం: లేదు. ఉంటే, పాక్షిక ప్రమేయం
\[
g(x) \simeq
\begin{cases}
0 & \text{$h(\gn{\cfind{x}(x)})=0$ అయితే} \\
\text{అనిర్వచితం} & \text{ఇతర సందర్భాల్లో}
\end{cases}
\]
గణనీయం అవుతుంది; కనుక దానికి ఏదో సూచిక~$e$ ఉంటుంది. కానీ అప్పుడు
\[
\cfind{e}(e) \simeq
\begin{cases}
0 & \text{$h(\gn{\cfind{e}(e)})=0$ అయితే} \\
\text{అనిర్వచితం} & \text{ఇతర సందర్భాల్లో,}
\end{cases}
\]
అవుతుంది. దీని ప్రకారం $\cfind{e}(e)$ నిర్వచితమైతేనే అది నిర్వచితం
కాదు; ఇది వైరుధ్యం.

ఇప్పుడు $\cfind{e}$ ఉన్న సమీకరణాన్ని చూడండి. ఒక అర్థంలో అక్కడ
స్వీయ-సూచన ఉంది: $\cfind{e}(e)$ విలువ ఒక నిర్దిష్ట విధంగా
$\gn{\cfind{e}(e)}$పై ఆధారపడేలా ఏర్పాటుచేశాం. వైరుధ్యాలను
నిరూపించడానికే కాకుండా, సాధారణంగా కూడా మనం ఇలా చేయగలమని స్థిరబిందు
సిద్ధాంతం చెబుతుంది.

స్థిరబిందు సిద్ధాంతాన్ని ప్రకటించడానికి రెండు తుల్య మార్గాలను
\olref{lem:fixed-equiv} ఇస్తుంది. తార్కికంగా చూస్తే రెండు ప్రకటనలూ
నిజాలు కనుక అవి తుల్యాలని వస్తుంది; కాని మన అసలు ఉద్దేశం, ప్రతి
ప్రకటనా మరొకదాని నుంచి నేరుగా అనుసరిస్తుంది కనుక అవి ఒకే సిద్ధాంతపు
ప్రత్యామ్నాయ ప్రకటనలుగా తీసుకోవచ్చని చెప్పడం.

\begin{lem}
\ollabel{lem:fixed-equiv}
కింది ప్రకటనలు తుల్యాలు:
\begin{enumerate}
\item ప్రతి పాక్షిక గణనీయ ప్రమేయం $g(x,y)$కు, ప్రతి~$y$కూ
\[
\cfind{e}(y) \simeq g(e,y).
\]
అయ్యే సూచిక~$e$ ఒకటి ఉంటుంది.
\item ప్రతి గణనీయ ప్రమేయం~$f(x)$కు, ప్రతి~$y$కూ
\[
\cfind{e}(y) \simeq \cfind{f(e)}(y).
\]
అయ్యే సూచిక~$e$ ఒకటి ఉంటుంది.
\end{enumerate}
\end{lem}

\begin{proof}
$(1)\Rightarrow(2)$: $f$ ఇచ్చినప్పుడు, $g(x,y)\simeq
\fn{Un}(f(x),y)$గా $g$ను నిర్వచించండి. ప్రతి~$y$కు
\begin{align*}
\cfind{e}(y) & = \fn{Un}(f(e),y) \\
& = \cfind{f(e)}(y).
\end{align*}
అయ్యే సూచిక~$e$ను పొందడానికి (1)ని వాడండి.

$(2)\Rightarrow(1)$: $g$ ఇచ్చినప్పుడు, ప్రతి $x$,~$y$లకు
$\cfind{f(x)}(y)\simeq g(x,y)$ అయ్యే $f$ను పొందడానికి $s$-$m$-$n$
సిద్ధాంతాన్ని వాడండి. ఆపై
\begin{align*}
\cfind{e}(y) & = \cfind{f(e)}(y) \\
& = g(e,y).
\end{align*}
అయ్యే సూచిక~$e$ను పొందడానికి (2)ని వాడండి. దీంతో నిరూపణ పూర్తయింది.
\end{proof}

\begin{explain}
ప్రకటన~(1) నిజమని (అందువల్ల (2) కూడా నిజమని) చూపడానికి ముందు, అది
ఎంత విచిత్రమైనదో పరిశీలించండి. $e$ను ఒక కంప్యూటర్ ప్రోగ్రాముగా
భావించండి. ఏ పాక్షిక గణనీయ $g(x,y)$ ఇచ్చినా, ఆ ప్రోగ్రామునే సూచించే
ప్రమేయం $g_e(y)\simeq g(e,y)$ను గణించే కంప్యూటర్ ప్రోగ్రాము~$e$ను
కనుగొనవచ్చని ప్రకటన~(1) చెబుతుంది. మరో మాటలో, తన ప్రోగ్రామునే
సూచించే ప్రమేయాన్ని గణించే కంప్యూటర్ ప్రోగ్రామును కనుగొనవచ్చు.
\end{explain}

\begin{thm}
\olref{lem:fixed-equiv}లోని రెండు ప్రకటనలూ నిజాలు. మరింత స్పష్టంగా,
ప్రతి పాక్షిక గణనీయ ప్రమేయం $g(x,y)$కు, ప్రతి~$y$కూ
\[
\cfind{e}(y) \simeq g(e,y).
\]
అయ్యే సూచిక~$e$ ఒకటి ఉంటుంది.
\end{thm}

\begin{proof}
కావలసిన అంశాలు పైన ఆగే సమస్య చర్చలో ఇప్పటికే అంతర్లీనంగా ఉన్నాయి.
ప్రతి~$x$కు $f_x(y)\simeq\cfind{x}(x,y)$ అనే ప్రమేయానికి ఒక సూచికను
తిరిగి ఇచ్చే గణనీయ ప్రమేయం~$\fn{diag}(x)$ను తీసుకుందాం; అంటే
\[
\cfind{\fn{diag}(x)}(y) \simeq \cfind{x}(x,y).
\]
2-స్థాన ప్రమేయానికి గల ప్రోగ్రామును, అసలు ప్రోగ్రామునే దాని మొదటి
ఆర్గుమెంటుగా స్థిరపరచి పొందిన 1-స్థాన ప్రమేయపు ప్రోగ్రాముగా మార్చే
ప్రమేయంగా $\fn{diag}$ను భావించండి. $\fn{diag}$ను కచ్చితంగా ఇలా
నిర్వచించవచ్చు: ముందుగా
\[
s(x,y)\simeq\fn{Un}^2(x,x,y),
\]
గా $s$ను నిర్వచించండి; ఇక్కడ $\fn{Un}^2$ అనేది పాక్షిక గణనీయ
2-స్థాన ప్రమేయాలకు సార్వత్రికమైన 3-స్థాన ప్రమేయం. అప్పుడు
$s$-$m$-$n$ సిద్ధాంతం ప్రకారం
\[
\cfind{\fn{diag}(x)}(y)\simeq s(x,y).
\]
నెరవేర్చే ఆదిమ పునరావృత్త ప్రమేయం $\fn{diag}$ను కనుగొనవచ్చు.

ఇప్పుడు ప్రమేయం~$l$ను
\[
l(x,y)\simeq g(\fn{diag}(x),y).
\]
అని నిర్వచించి, $\gn{l}$ అనేది $l$కు ఒక సూచిక అనుకుందాం. చివరగా
$e=\fn{diag}(\gn{l})$ అని ఉంచండి. అప్పుడు ప్రతి~$y$కు
\begin{align*}
\cfind{e}(y) & \simeq \cfind{\fn{diag}(\gn{l})}(y) \\
& \simeq \cfind{\gn{l}}(\gn{l}, y) \\
& \simeq l(\gn{l}, y) \\
& \simeq g(\fn{diag}(\gn{l}),y) \\
& \simeq g(e, y),
\end{align*}
కావలసినట్లే.
\end{proof}

\begin{explain}
ఇక్కడ జరుగుతున్నది ఏమిటి? తననే తాను ముద్రించే కంప్యూటర్ ప్రోగ్రామును
రాయాలని మీకు పని ఇచ్చారనుకోండి. అయితే మీరు అక్షరమాల ప్రమేయాల విచిత్రమైన,
సమృద్ధమైన గ్రంథాలయం గల ప్రోగ్రామింగ్ భాషను వాడుతున్నారనుకుందాం.
ప్రత్యేకంగా, మీ భాషలో ఇలా పనిచేసే ప్రమేయం~$\fn{diag}$ ఉందనుకోండి:
ఇన్‌పుట్ అక్షరమాల~$s$ ఇచ్చినప్పుడు, $s$లో కనిపించే `x' సంకేతపు ప్రతి
సందర్భాన్నీ $\fn{diag}$ గుర్తించి, అసలు అక్షరమాల యొక్క ఉల్లేఖిత రూపంతో
దాన్ని భర్తీ చేస్తుంది. ఉదాహరణకు, ఇన్‌పుట్‌గా
\begin{quote}
\begin{verbatim}
hello x world
\end{verbatim}
\end{quote}
అనే అక్షరమాలను ఇస్తే, ఆ ప్రమేయం
\begin{quote}
\begin{verbatim}
hello 'hello x world' world
\end{verbatim}
\end{quote}
అనే అవుట్‌పుట్‌ను ఇస్తుంది. అప్పుడు కావలసిన ప్రోగ్రామును రాయడం సులభం;
\begin{quote}
\begin{verbatim}
print(diag('print(diag(x))'))
\end{verbatim}
\end{quote}
పని చేస్తుందని తనిఖీ చేయవచ్చు. C++, Java వంటి మరింత సాధారణ
ప్రోగ్రామింగ్ భాషలకు కూడా ఇదే భావన (మరింత క్లిష్టమైన అమలుతో) పనిచేస్తుంది.

స్థిరబిందు సిద్ధాంతపు నిరూపణకు మనం ఇంకొన్ని అడుగుల దూరంలోనే ఉన్నాం.
$\fn{print}(x,y)$ అనే ముద్రణ ప్రమేయపు ఒక రూపాంతరం అక్షరమాల~$x$ను,
మరొక సంఖ్యాత్మక ఆర్గుమెంటు~$y$ను తీసుకుని, $x$ను $y$ సార్లు
ముద్రిస్తుందనుకోండి. అప్పుడు ``ప్రోగ్రాము''
\begin{quote}
\begin{verbatim}
getinput(y);
print(diag('getinput(y); print(diag(x), y)'), y)
\end{verbatim}
\end{quote}
ఇన్‌పుట్~$y$పై తననే $y$ సార్లు ముద్రిస్తుంది.
$\fn{getinput}$---$\fn{print}$---$\fn{diag}$ అస్థిపంజరాన్ని యాదృచ్ఛిక
ప్రమేయం $g(x,y)$తో భర్తీ చేస్తే
\begin{quote}
\begin{verbatim}
g(diag('g(diag(x), y)'), y)
\end{verbatim}
\end{quote}
వస్తుంది. ఇది ఇన్‌పుట్~$y$పై $g$ను ఆ ప్రోగ్రాము మీదనే, అలాగే $y$ మీద
నడిపే ప్రోగ్రాము. ``ఉల్లేఖించడం'' అంటే ``దానికి ఒక సూచికను వాడటం''గా
భావిస్తే, పై నిరూపణ లభిస్తుంది.

ప్రస్తుతానికి ఈ నిరూపణను కేవలం ఆచారిక చాకచక్యంగా లేదా మాయగా
భావించాలనుకుంటే పరవాలేదు. కానీ పైన ఇచ్చిన వాదన వివరాలను మీరు తిరిగి
నిర్మించగలగాలి. అసంపూర్ణతా సిద్ధాంతాలను (మరియు సంబంధిత ``స్థిరబిందు
సిద్ధాంతాన్ని'') నిరూపించేటప్పుడు ఇది ఎందుకు పనిచేస్తుందో అర్థం చేసుకునే
ఇతర మార్గాలను చర్చిస్తాం.
\end{explain}

\begin{tagblock}{lambda}
\begin{digress}
ఇదే భావనతో ఒక ``స్థిరబిందు'' సంయోజకాన్ని పొందవచ్చు. మీ దగ్గర లాంబ్డా
పదం~$g$ ఉందనీ, మరొక పదం~$k$ కావాలనీ, ఆ $k$ అనేది $gk$కు
$\beta$-తుల్యం కావాలనీ అనుకుందాం. మన సంకేత సంప్రదాయాలతో
\[
\fn{diag}(x)=xx
\]
మరియు
\[
l(x)=g(\fn{diag}(x))
\]
అనే పదాలను నిర్వచించండి; మరో మాటలో, $l$ అనేది
$\lambd[x][g(xx)]$ అనే పదం. $k$ను $ll$ అనే పదంగా ఉంచండి. అప్పుడు
\begin{align*}
k & = (\lambd[x][g(xx)])(\lambd[x][g(xx)]) \\
& \red g((\lambd[x][g(xx)])(\lambd[x][g(xx)])) \\
& = gk.
\end{align*}
\[
Y=\lambd[g][((\lambd[x][g(xx)])(\lambd[x][g(xx)]))]
\]
అని తీసుకుంటే $Yg$, $g(Yg)$ రెండూ ఒకే పదానికి తగ్గుతాయి; కనుక
$Yg\equiv_\beta g(Yg)$. దీన్ని ``కరీ సంయోజకం'' అంటారు. బదులుగా
\[
Y=(\lambd[xg][g(xxg)])(\lambd[xg][g(xxg)])
\]
అని తీసుకుంటే, నిజానికి $Yg$ అనేది $g(Yg)$కు తగ్గుతుంది; ఇది మరింత
బలమైన ప్రకటన. $Y$ యొక్క ఈ రెండవ రూపాన్ని ``ట్యూరింగ్ సంయోజకం'' అంటారు.
\end{digress}
\end{tagblock}

\end{document}
